\section{Error Analysis and Limitations}

We manually inspect 100 development examples for which CEA and mean pooling make
different predictions. In 61 cases, CEA assigns less weight to a sentence that shares
keywords with the claim but does not report the relevant experiment. Another 24 cases
contain conflicting measurements from two abstracts; confidence calibration helps when
one source states its experimental conditions more clearly. The remaining examples are
split between annotation ambiguity and retrieval failures.

CEA cannot recover evidence that is absent from the retrieved set. It also treats
source confidence as a property learned from textual representations rather than from
external information about venues, authors, or replication history. Consequently, a
fluent but incorrect statement may still receive a high confidence score. These limits
suggest that calibrated aggregation complements, but does not replace, careful evidence
retrieval and provenance tracking.

EvidenceBench is synthetic and substantially smaller than a production scientific
corpus. The reported numbers should therefore be interpreted as controlled measurements
of model behavior rather than estimates of real-world fact-checking performance. We do
not evaluate multilingual claims, full-document retrieval, or evidence expressed only
in figures. Extending the evaluation along these dimensions is necessary before using
the method in a high-stakes review workflow.
